\documentclass[12pt,a4paper]{article}
\usepackage[utf8]{inputenc}
\usepackage[T1]{fontenc}
\usepackage[romanian]{babel}
\usepackage{graphicx}
\usepackage{hyperref}
\usepackage{setspace}
\usepackage{geometry}

\geometry{margin=2.5cm}
\onehalfspacing

\title{\textbf{Abacul roman}}
\author{
Student 1: NUME PRENUME \\
Student 2: NUME PRENUME
}
\date{}

\begin{document}

\maketitle
\thispagestyle{empty}

\newpage

\section*{Introducere}

Încă din cele mai vechi timpuri, oamenii au avut nevoie de metode eficiente pentru efectuarea calculelor matematice. Activități precum comerțul, construcțiile, administrarea bunurilor sau colectarea taxelor necesitau calcule rapide și corecte. Înainte de apariția sistemului modern de numerotație și a calculatoarelor electronice, civilizațiile antice au dezvoltat diferite instrumente pentru a simplifica operațiile aritmetice. Unul dintre cele mai importante astfel de instrumente a fost abacul.

Abacul roman reprezintă un dispozitiv mecanic de calcul utilizat în perioada Imperiului Roman. Acesta poate fi considerat unul dintre cele mai vechi „calculatoare” portabile din istorie. Dispozitivul era folosit pentru efectuarea rapidă a operațiilor aritmetice, cum ar fi adunarea, scăderea, înmulțirea și împărțirea. Prin intermediul său, comercianții, inginerii și funcționarii romani puteau realiza calcule mult mai eficient decât folosind doar numeralele romane. \cite{romanabacuswiki}

Importanța abacului roman este dată atât de rolul său practic în viața economică și administrativă a Imperiului Roman, cât și de contribuția sa la evoluția instrumentelor de calcul. Deși sistemul de numerotație roman nu includea cifra zero și era relativ dificil de utilizat pentru calcule complexe, abacul permitea efectuarea operațiilor matematice într-un mod mai rapid și mai intuitiv. \cite{ancientcomputers}

\section*{Originea și evoluția abacului}

Cuvântul „abac” provine din termenul grecesc „abax”, care înseamnă „placă” sau „tablă”. Primele forme de abac erau foarte simple și constau în suprafețe pe care erau plasate pietricele sau alte obiecte mici pentru a reprezenta numere. Aceste instrumente au apărut în diferite regiuni ale lumii, inclusiv în Mesopotamia, Grecia și China. \cite{unrvabacus}

Romanii au preluat ideea acestor instrumente de calcul și au dezvoltat propriul lor tip de abac, adaptat nevoilor societății romane. Inițial, calculul era realizat pe table de numărare, unde erau plasate mici pietre sau jetoane pentru a reprezenta diferite valori numerice. În timp, aceste dispozitive au evoluat într-o formă mai compactă și mai practică, cunoscută sub numele de abacul roman portabil. \cite{romanhandabacus}

În jurul primului secol al erei noastre, romanii au început să producă abacuri realizate din metal, cu șanțuri în care se deplasau bile sau mici discuri. Aceste dispozitive aveau dimensiuni reduse și puteau fi transportate cu ușurință, fiind utilizate de comercianți și funcționari în activitățile zilnice. Din acest motiv, abacul roman este considerat unul dintre primele instrumente de calcul portabile din istorie. \cite{romanabacuswiki}

\section*{Structura și modul de funcționare}

Abacul roman era format, de obicei, dintr-o placă metalică prevăzută cu mai multe șanțuri paralele. În interiorul acestor șanțuri se aflau mici bile sau discuri care puteau fi deplasate pentru a reprezenta diferite valori numerice. Fiecare coloană a abacului corespundea unei anumite unități de valoare, cum ar fi unități, zeci, sute sau mii. \cite{romanhandabacus}

În majoritatea coloanelor existau două tipuri de bile: cele din partea inferioară și cele din partea superioară. Bilele din partea inferioară reprezentau unități simple, fiecare având valoarea unu, în timp ce bila din partea superioară avea valoarea cinci. Prin combinarea acestor valori, utilizatorul putea reprezenta orice cifră de la zero la nouă într-o coloană. Acest sistem este cunoscut sub numele de sistem bi-cinar și este similar cu modul de funcționare al unor tipuri moderne de abac. \cite{ancientcomputers}

Un aspect interesant al abacului roman este faptul că unele coloane erau dedicate fracțiilor. Romanii foloseau frecvent sistemul duodecimal, bazat pe împărțirea unei unități în douăsprezece părți. De exemplu, moneda romană principală, numită „as”, era împărțită în douăsprezece unități mai mici numite „unciae”. Din acest motiv, abacul roman includea coloane speciale pentru calculul acestor fracții. \cite{romanabacuswiki}

Prin deplasarea bilelor de-a lungul șanțurilor, utilizatorul putea efectua operații matematice relativ complexe. Calculul era realizat prin manipularea directă a acestor bile, iar rezultatul final putea fi apoi transcris folosind numeralele romane. \cite{romanhandabacus}

\section*{Utilizarea abacului în societatea romană}

Abacul roman a avut un rol important în viața economică și administrativă a Imperiului Roman. Comerțul intens, taxele și proiectele de construcție necesitau calcule precise și rapide. În acest context, abacul a devenit un instrument esențial pentru comercianți, contabili, ingineri și colectori de taxe. \cite{unrvabacus}

De exemplu, comercianții foloseau abacul pentru a calcula prețuri, cantități și taxe comerciale. În același timp, inginerii romani îl utilizau pentru a efectua calcule necesare în construcția drumurilor, podurilor sau apeductelor. Deși rezultatele erau în final scrise folosind numeralele romane, calculele propriu-zise erau realizate pe abac, deoarece sistemul de numerotație roman nu era potrivit pentru operații aritmetice complexe. \cite{romanhandabacus}

Abacul era utilizat și în educație, fiind un instrument important pentru învățarea matematicii. Copiii puteau exersa operațiile de bază folosind acest dispozitiv, ceea ce contribuia la dezvoltarea abilităților de calcul. \cite{unrvabacus}

\section*{Importanța istorică a abacului roman}

Deși astăzi este considerat un instrument simplu, abacul roman reprezintă un pas important în evoluția tehnologiilor de calcul. Acesta a demonstrat că manipularea obiectelor fizice poate fi utilizată pentru a reprezenta și procesa informații numerice.

Abacul roman a influențat dezvoltarea altor tipuri de abac, precum abacul chinezesc (suanpan) sau cel japonez (soroban). Deși aceste instrumente au evoluat independent și au fost adaptate diferitelor sisteme numerice, principiul de bază a rămas același: utilizarea unor bile sau marcatori pentru reprezentarea valorilor numerice și efectuarea calculelor. \cite{ancientcomputers}

Mai mult decât atât, abacul poate fi considerat un precursor al calculatoarelor moderne. Chiar dacă funcționarea sa era complet mecanică și manuală, el realiza aceleași operații fundamentale care stau la baza calculatoarelor actuale: reprezentarea numerelor și manipularea lor pentru a obține rezultate. \cite{romanabacuswiki}

\section*{Concluzie}

Abacul roman reprezintă una dintre cele mai importante invenții din domeniul instrumentelor de calcul din Antichitate. Prin designul său simplu, dar eficient, acesta a permis realizarea rapidă a calculelor într-o perioadă în care sistemele numerice erau limitate și dificil de utilizat.

Rolul său în societatea romană a fost semnificativ, fiind utilizat în comerț, administrație, inginerie și educație. De asemenea, abacul roman a contribuit la dezvoltarea ulterioară a altor instrumente de calcul și poate fi considerat un precursor al tehnologiilor moderne de calcul.

Chiar dacă astăzi calculatoarele electronice au înlocuit complet aceste dispozitive simple, abacul rămâne un simbol important al ingeniozității umane și al dorinței de a găsi metode eficiente pentru rezolvarea problemelor matematice.

\newpage

\begin{thebibliography}{9}

\bibitem{romanhandabacus}
Steve Stephenson, \textit{The Roman Hand-Abacus}.  
Disponibil la: \url{https://www.ee.torontomu.ca/~elf/abacus/roman-hand-abacus.html}

\bibitem{ancientcomputers}
Engineering and Technology History Wiki, \textit{Ancient Computers}.  
Disponibil la: \url{https://ethw.org/Ancient_Computers}

\bibitem{classroomadventures}
Classroom Adventures, \textit{The Roman Abacus}.  
Disponibil la: \url{https://www.classroomadventures.co.uk/post/the-roman-abacus}

\bibitem{romanabacuswiki}
Wikipedia, \textit{Roman abacus}.  
Disponibil la: \url{https://en.wikipedia.org/wiki/Roman_abacus}

\bibitem{unrvabacus}
UNRV Roman History, \textit{Roman Abacus}.  
Disponibil la: \url{https://www.unrv.com/culture/abacus.php}

\end{thebibliography}

\end{document}